\subsection{KT-ARTIN-FAIL: an affine tower of elementary fibrations}
\label{sec:kt-artin-fail}

\begin{proposition}[KT-ARTIN-FAIL]
The counterexample to prime-to-$p$ point separation can be taken
affine, with both steps in its polycurve tower elementary
fibrations in hyperbolic curves. It therefore applies to Artin
neighbourhoods understood to mean such towers, without the
additional product-embedding condition of a \emph{strongly}
hyperbolic Artin neighbourhood.
\end{proposition}

\begin{proof}
Keep $k$, $B$, $B'$, $\omega$ and $b=\omega$ from (PF.1).
Replace $F$ by the smooth projective curve
\[
 F:\quad y^2+y=x+\frac1x+\frac1{x+1}.
                                                               \tag{AF.1}
\]
Its three simple poles are $0,1,\infty$, so it has genus two
by Stichtenoth, Proposition~3.7.8(c)--(d). The corresponding
three ramification points $q_0,q_1,q_\infty$ are $k$-rational
and fixed by the involution $\iota:y\mapsto y+1$.
Put $F^\circ=F\setminus\{q_\infty\}$ and
\[
 X^\circ=(F^\circ\times B')/\langle(\iota,\tau)\rangle.
                                                               \tag{AF.2}
\]
This is a smooth affine surface. Its morphism to $B$ has a
smooth proper compactification obtained from $F\times B'$,
with complement the single section descended from $q_\infty$.
Both $X^\circ\to B$ and $B\to\operatorname{Spec}k$ are
therefore elementary fibrations with hyperbolic fibres.
The point $q_0$ gives a section of $X^\circ\to B$.

Fix compatible geometric base points using this section.
Let $\varphi:\pi_1(X^\circ_{\overline k})\twoheadrightarrow Q$
be a finite quotient of odd order. The image of the index-two
subgroup
\[
 \pi_1(F^\circ_{\overline k}\times B'_{\overline k})
 \subset\pi_1(X^\circ_{\overline k})
\]
has index dividing both two and $|Q|$, hence is all of $Q$.
Prime-to-two K\"unneth, SGA~1, Expos\'e XIII,
Proposition~4.6, pp.~310--311, applies to these smooth curves:
the required strong resolution is only in dimension at most
one. Consequently $Q$ is generated by the commuting images
$A$ of $\pi_1(F^\circ_{\overline k},q_0)$ and $H'$ of
$\pi_1(B'_{\overline k})$. In particular $A$ is normal in $Q$.

Take an element of $\pi_1(B_{\overline k})$ with nontrivial
double-cover character, and lift it using the $q_0$ section.
In this finite odd quotient, conjugation by its image on $A$
is induced by $\iota_*$. Indeed, the product trivialization
over $B'$ identifies the descent action on the pointed fibre
with $\iota$; any base-cover transport acts through $H'$, which
commutes with $A$ by prime-to-two K\"unneth. Only this relation
in $Q$ is used, not an exact full wild homotopy sequence for
the affine fibration. Consequently that conjugation on $A$
has order dividing two. It also has odd order, since $Q$ has
odd order. Thus $\iota_*$ acts trivially on $A$.

Let $Y\to F^\circ_{\overline k}$ be the connected \'etale
Galois cover with group $A$. The pointed trivial action of
$\iota_*$ gives a lift $j$ of $\iota$ of order two, fixing a
chosen point above $q_0$ and commuting with $A$.
The quotient map $F^\circ\to\mathbf A^1_x$ is finite, so
the normalization of $\mathbf A^1$ in the function field of
$Y$ is $Y$ itself. The composite $Y\to\mathbf A^1$ is
Galois with group $A\times\langle j\rangle$. Its inertia
at every point of $\mathbf A^1$ is either trivial or of
order two, because $Y\to F^\circ$ is \'etale. The only
order-two element in the direct product is $j$.
It follows that
\[
 Y/\langle j\rangle\longrightarrow\mathbf A^1_{\overline k}
\]
is \'etale and Galois of odd degree $|A|$.

There is no nontrivial finite odd Galois \'etale cover of the
geometric affine line in characteristic two. Indeed, its
smooth projective compactification can ramify only at infinity,
and ramification is tame because its Galois group has odd order.
If its degree is $d>1$ and the ramification index at infinity
is $e$, tame Riemann--Hurwitz reads
\[
 2g-2=-2d+\frac d e(e-1)=-d-\frac d e<-2,
\]
a contradiction. Hence $A=1$.

Since this holds for every finite odd quotient, the geometric
fibre group maps trivially to
$\pi_1(X^\circ_{\overline k})^{(2')}$.
The two distinct points of the fibre over $b$ corresponding
to $q_0$ and $q_1$ are rational over $k$ and have equal
arithmetic section classes in $\Pi_{X^\circ}^{(2')}$,
by the last argument of KT-POLY-FAIL.
\end{proof}

The result is compatible with KT-PRODUCT-SEP: an embedding
in a product of hyperbolic curves would contradict the
exhibited section collision, so this surface does not meet
that additional condition.

